\section{Frecuencia de muestreo de un sistema de control continuo}
\label{sec:p2}
\enunciado{¿Cómo determinar la frecuencia de muestreo de un sistema de control continuo? Proponga
una frecuencia razonable para el proyecto suyo (15\,\%).}

\subsection{Criterios}
Un controlador digital observa la planta solo en los instantes $t_k=kT_s$ y mantiene su salida
constante entre ellos con un retenedor de orden cero (ZOH). El teorema de Nyquist garantiza que una
señal de banda limitada se pueda reconstruir, pero en control no basta: el ZOH equivale
aproximadamente a un retardo de $T_s/2$, que resta al margen de fase
\begin{equation}
  \Delta\varphi \approx \frac{\omega_c T_s}{2}\ \text{[rad]},
  \label{eq:fase_zoh}
\end{equation}
donde $\omega_c$ es la frecuencia de cruce del lazo~\cite{franklin1998digital,astrom1997computer}.
Además, las aproximaciones de $s$ que convierten el PID en una ecuación en diferencias solo son
fieles si $T_s$ es pequeño frente a las constantes de tiempo del lazo~\cite{ogata1996discreto}.
Por eso se usan reglas prácticas: muestrear de 10 a 40 veces el ancho de banda de lazo cerrado
($\omega_s\approx 20\,\omega_b$)~\cite{franklin1998digital}, mantener $\omega_c T_s$ entre 0,15 y
0,5, y tener de 4 a 10 muestras por tiempo de subida~\cite{astrom1997computer}. El periodo también
debe ser compatible con el actuador, el sensor y el tiempo de cálculo.

\subsection{Propuesta: $T_s=\SI{20}{ms}$ ($f_s=\SI{50}{Hz}$)}
Para un frenado de 1,5 a 4~s y un servo que tarda unos 0,1~s en girar 60°~\cite{mg90s}, un ancho
de banda de lazo cerrado razonable es $\omega_b\approx\omega_c\approx$ 4 a $\SI{5}{rad/s}$. La
Tabla~\ref{tab:ts} verifica el periodo propuesto.

\begin{table}[h]
\centering\small
\caption{Verificación de $T_s=\SI{20}{ms}$ con $\omega_c=\SI{4}{rad/s}$.}
\label{tab:ts}
\begin{tabularx}{\linewidth}{@{}>{\raggedright\arraybackslash}p{4cm}>{\raggedright\arraybackslash}X>{\raggedright\arraybackslash}p{3.4cm}@{}}
\toprule
Criterio & Cálculo & Resultado \\
\midrule
$\omega_s \gtrsim 20\,\omega_b$ & $\omega_s=2\pi/T_s=\SI{314}{rad/s}$ frente a $\SI{100}{rad/s}$ & Cumple \\
$\omega_c T_s \leq 0{,}5$ & $4\cdot 0{,}02 = 0{,}08$ & Cumple \\
Fase del ZOH, ec.~\eqref{eq:fase_zoh} & $0{,}08/2=\SI{0{,}04}{rad}$ & $2{,}3^\circ$, despreciable \\
Muestras por subida & $t_r\approx 2{,}2/\omega_b\approx\SI{0{,}5}{s}$; $0{,}5/0{,}02$ & 25 \\
Actuador & El servo lee un pulso cada 20~ms~\cite{mg90s} & Muestrear más rápido no se aplica \\
\bottomrule
\end{tabularx}
\end{table}

El periodo de control no depende de la velocidad máxima del disco, que solo sirve para validar el
sensor. La velocidad se mide por tiempo entre marcas y no contando marcas en una
ventana~\cite{petrella2007speed}: con $N=10$, contar en 20~ms daría una resolución de
$2\pi/(N T_s)=\SI{31{,}4}{rad/s}$, mientras que medir el intervalo con resolución de 1~µs da
$0{,}024\,\%$ a $\SI{150}{rad/s}$.

Se propone $T_s=\SI{20}{ms}$: cumple las reglas con margen y coincide con la trama del servo. La
caracterización de la pregunta~7 lo confirma: la constante de tiempo medida, $\tau=\SI{0{,}274}{s}$,
corresponde a unas 14 muestras por constante de tiempo.
